% ==================================================================
% id:       lim-042
% topic:    גבולות ורציפות של פונקציות
% subtopic: אינסוף כגבול וגבול באינסוף
% solution: full
% ==================================================================
מצאו $a$ כך ש-$\ds\lim_{x\to\infty}\lr{\sqrt{x^2+ax+1}-x}=2$.
\ifsol
\begin{solbox}
\[\sqrt{x^2+ax+1}-x=\frac{ax+1}{\sqrt{x^2+ax+1}+x}=\frac{a+\frac1x}{\sqrt{1+\frac ax+\frac1{x^2}}+1}\]
ולכן
\[\lim_{x\to\infty}\left(\sqrt{x^2+ax+1}-x\right)=\lim_{x\to\infty}\left(\frac{a+\frac1x}{\sqrt{1+\frac ax+\frac1{x^2}}+1}\right)=\frac a2\]
לכן $\frac a2=2$, כלומר $a=4$.
הערה: בחישוב הגבול השתמשנו באריתמטיקה של גבולות: הגבול של סכום, של מכפלה ושל מנה של פונקציות מתכנסות שווה לסכום, למכפלה ולמנה של הגבולות, כל עוד הגבול של המכנה שונה מ-$0$ בנוסף, השורש הוא פונקציה רציפה, ולכן אפשר להעביר את הגבול אל תוך השורש.
\end{solbox}
\fi
